\chapter{Introducción}
\label{cap:introduccion}
\section{Contexto}
Este Trabajo Final de Ingeniería en Informática aborda el diseño y la evaluación de
Neuma, una plataforma de entrenamiento auditivo musical para dispositivos Android.
La propuesta reúne ejercicios presentados como juegos serios, análisis de respuestas
vocales y ajuste de la dificultad según el desempeño del usuario. El proyecto se
encuadra en la modalidad de Proyecto y/o Diseño de Ingeniería en Informática.

La aplicación está dirigida principalmente a estudiantes de conservatorio y, de manera
más amplia, a estudiantes de música e interesados en ejercitar su oído y su lectura
musical. Se presupone una base de conocimientos musicales, incluida la lectura de
partituras necesaria para las actividades que la utilizan. El público previsto comprende
niños, adolescentes y adultos. Neuma se concibe como un recurso complementario de las
clases o del aprendizaje realizado por otros medios.

El escenario de uso considerado es la práctica individual: el usuario escucha una
consigna musical, responde mediante los controles del juego o mediante la voz y recibe
información sobre su respuesta. Según la actividad, el sistema debe interpretar una
secuencia de notas, una elección tímbrica, una ejecución rítmica o una respuesta vocal.
Estas interacciones plantean necesidades diferentes de reproducción, captura, evaluación
y presentación de la información, que deben integrarse en una misma aplicación.

En los ejercicios vocales, la respuesta requiere transformar la señal del micrófono en
una estimación de frecuencia fundamental y relacionarla con el objetivo musical. En los
demás ejercicios, el desempeño se obtiene de las acciones registradas durante la partida.
En ambos casos, la propuesta contempla utilizar indicadores de desempeño para regular
los desafíos y conservar información que permita examinar las decisiones del sistema.
El marco teórico y los antecedentes que fundamentan estos componentes se desarrollarán
en el capítulo~\ref{cap:marco}.
\section{Motivación}
El proyecto surge de mi experiencia como estudiante de música en el
conservatorio. Durante mi formación encontré dificultades para practicar determinados
contenidos y limitaciones en las herramientas que utilizaba: algunas no se ajustaban
a mi forma de estudiar y otras restringían el acceso mediante planes pagos. También
percibí que, en ciertas actividades, mejorar el desempeño dentro del juego no tenía
una relación clara con la habilidad musical que deseaba desarrollar. En conversaciones
informales con otros estudiantes de música de la ciudad reconocí
inquietudes similares. Estas experiencias constituyen el origen de la propuesta y
no representan un relevamiento sistemático de necesidades ni una evaluación comparativa
de las aplicaciones disponibles.

La motivación técnica del trabajo es construir una integración verificable entre
procesamiento de audio, interacción musical y adaptación de ejercicios. Mostrar una nota
estimada o modificar un parámetro de dificultad constituye una parte de esa integración;
también es necesario establecer bajo qué condiciones funciona, qué información utiliza
y cuáles son sus límites. Por ello, el desarrollo del prototipo se acompaña de un registro
de datos y de una evaluación de su comportamiento.

Desde la perspectiva de la práctica musical, busco ofrecer actividades en las que el
usuario pueda reconocer el resultado de sus respuestas y continuar ejercitándose con
desafíos ajustados a su desempeño. Esta finalidad orienta el diseño de la aplicación.
Para orientar las actividades hacia ese propósito, el diseño deberá explicitar qué
habilidad musical practica cada una e incorporar variaciones pertinentes que reduzcan
la dependencia de pistas particulares del minijuego. Por ejemplo, el ajuste de afinación
de oído busca ejercitar la escucha de diferencias de altura, mientras que la variación
tímbrica en ejercicios de intervalos pretende evitar que la respuesta dependa siempre
de una misma sonoridad. La pertinencia de cada variación deberá examinarse según la
actividad. Adopto la transferencia a situaciones musicales externas a la aplicación
como intención de diseño; demostrarla requeriría una evaluación específica que excede
las pruebas técnicas y exploratorias previstas en este trabajo.

El interés académico reside en documentar las decisiones de ingeniería que permiten
construir el prototipo y contrastarlas con evidencia obtenida en condiciones definidas.
El aporte propuesto comprende el software, su diseño y el análisis de precisión de la
estimación de tono, latencia de la retroalimentación y comportamiento del ajuste de
dificultad. La publicación gratuita del código y su documentación forma parte del
producto previsto en el protocolo.
\section{Problema a resolver}
El problema consiste en diseñar e integrar un sistema móvil que interprete las respuestas
del usuario, produzca retroalimentación durante la práctica y ajuste los ejercicios
de acuerdo con indicadores de desempeño. El protocolo presentado el 28 de septiembre
de 2026 plantea evaluar esa integración mediante la precisión de detección de tono,
la latencia de la retroalimentación y el funcionamiento del ajuste de dificultad.

La entrada vocal introduce un problema de medición: se debe obtener una estimación de
tono cuya exactitud y tiempo de disponibilidad puedan contrastarse con referencias.
El componente adaptativo introduce un problema de decisión: es necesario definir qué
se considera evidencia de desempeño, cómo se transforma en un ajuste y en qué momento
se aplica. La integración exige, además, que esas decisiones sean compatibles con las
reglas de cada actividad y puedan reconstruirse a partir de los registros del prototipo.

La pregunta que organiza el desarrollo es cómo integrar el análisis vocal, la
retroalimentación y la dificultad adaptativa en una aplicación Android, y cómo evaluar
su funcionamiento con criterios técnicos explícitos. Para responderla, se deberán
identificar las condiciones de ensayo, los dispositivos utilizados, las referencias de
comparación y los indicadores de evaluación. Los objetivos del
capítulo~\ref{cap:objetivos} descomponen este problema en resultados verificables.
\section{Alcance y delimitaciones}
El producto previsto es un prototipo funcional que integre captura y análisis de audio,
ejercicios musicales, retroalimentación, ajuste de dificultad y registro del desempeño.
El protocolo propone inicialmente cinco actividades: Simón musical armónico,
reconocimiento de instrumentos, defensa de acordes en partitura, eco rítmico y reproducción
vocal de intervalos. El plan de trabajo adjunto acota el mínimo a un juego vocal y otro
juego sencillo, con tres actividades adicionales opcionales. La selección inicial y el
mínimo planificado se conservan como referencias distintas para revisar el alcance final.

Como alcance de entrega propuesto, el 4 de noviembre de 2026 definí cuatro
actividades: intervalos vocales, reconocimiento de instrumentos, Rhythm Tap y
Defender. Incluyo la respuesta en piano virtual de Defender y la selección
automática de ayuda vocal como compromisos propuestos todavía pendientes de
desarrollo. Los requerimientos se reúnen en el capítulo~\ref{cap:requerimientos};
la aprobación académica de este alcance queda sujeta a revisión del director.
Simón permanece pendiente de reformulación, reemplazo o eliminación, y el
afinador se conserva como módulo complementario.

Se mantiene como compromiso mínimo la integración de un ejercicio de reproducción vocal
de intervalos: el usuario escuchará una referencia, responderá con la voz y recibirá
retroalimentación, con dificultad ajustada a su desempeño. Este ejercicio deberá
demostrar la articulación entre análisis vocal y adaptación. La prueba técnica de voz
y el afinador aportan componentes al desarrollo, pero no sustituyen por sí solos esa
integración.

Las decisiones de desarrollo posteriores fijan Android como plataforma, Kotlin y Jetpack
Compose para la aplicación, funcionamiento sin conexión y un perfil local. También
establecen un ajuste de dificultad mediante lógica difusa y una primera prueba de voz
separada, basada en YIN. Estas decisiones orientan la construcción incremental; la
existencia de esa prueba separada no acredita la integración del juego vocal previsto.

Existen diferencias técnicas entre los documentos de partida. El análisis de factibilidad
del protocolo incluye una capa C++ con Oboe, transporte mediante estructuras sin bloqueos
y un detector híbrido DSP/ML dentro del producto mínimo. En cambio, el plan de trabajo
considera Oboe y el aprendizaje automático como opciones por validar y condiciona este
último a los resultados de la línea base y al tiempo disponible. Para el desarrollo
adopto su tratamiento como alternativas cuya incorporación dependerá de las mediciones.
Esta decisión deberá justificarse mediante evidencia y
acordarse con el director por su diferencia con la formulación del protocolo.

La versión de entrega se prevé gratuita en todas sus funciones, sin suscripción y
utilizable sin conexión. Se contempla publicar el código fuente y la documentación
bajo una licencia abierta compatible con las dependencias y los recursos empleados.
La enseñanza de contenidos musicales desde sus fundamentos queda fuera del alcance
actual. Una eventual sección de aprendizaje introductorio se considera una ampliación
futura.

Se mantiene la delimitación a entrada monofónica y procesamiento en el dispositivo.
Quedan fuera la extracción de voz en entornos polifónicos, el multijugador en tiempo real
y la realidad extendida con motores tridimensionales pesados. La evaluación prevista
comprende pruebas técnicas y sesiones exploratorias de uso, inicialmente con estudiantes
de música mayores de edad que posean conocimientos musicales básicos. Esta delimitación
de los participantes es más acotada que el público previsto para la aplicación. Las
observaciones se interpretarán dentro de las condiciones y la población efectivamente
evaluadas. La realización de las sesiones y sus resultados se documentarán cuando
exista evidencia; no se plantean como demostración de eficacia pedagógica.
